\section{Results}
\label{sec:results}

\subsection{Item-level availability across 20 comparisons}
\label{sec:results:availability}

The 20 eligible comparisons require 206 evidence items between them, derived from each
comparison's protocol-component set and frozen before any retrieval. Table~\ref{tab:states}
gives the assigned states.

\begin{table}[t]
\caption{Item states across 206 required evidence items, and aggregate class across the 20
comparisons. No item reaches \textsc{public-hash-verified}: not one of the four studies
records an immutable digest for any run artifact, so retrieved bytes have nothing to verify
against.}
\label{tab:states}
\begin{tabular}{lr}
\toprule
Item state & Items \\
\midrule
\textsc{public-unverified} & 17 \\
\textsc{inaccessible} & 189 \\
\textsc{public-hash-verified} & 0 \\
\textsc{private} & 0 \\
\textsc{reported-not-released} & 0 \\
\textsc{not-reported} & 0 \\
\midrule
Aggregate class & Comparisons \\
\midrule
\textsc{public-partial} & 10 \\
\textsc{private-or-unavailable} & 10 \\
\textsc{public-complete} & 0 \\
\bottomrule
\end{tabular}
\end{table}

Two of the zeros carry different weight. \textsc{public-hash-verified} is unreachable
because no source records a digest to verify against, which is a property of the sources
rather than of the retrieval. \textsc{not-reported} is empty because every comparison has at
least one required item that some attempted locator did carry, so incomplete closures route
to \textsc{inaccessible} instead.

Counting what blocks the 189 inaccessible items by kind, the largest is terminal and
missingness evidence at 134, followed by evaluator outcomes at 69 and analysis binding at
62. The first two are absences of data. The third is not: those closures fail because no
source states which released artifact produced which reported number.

\subsection{Bounded checks of released evidence}
\label{sec:results:recompute}

The operations in this subsection differ by study and the claims are limited accordingly;
Table~\ref{tab:evidence-operation-claim} states the ceiling each one carries. For
the selected SWE-Gym~\cite{TrainingSoftwareEngineering2024} Moatless subset, we count the
released \emph{resolved} and \emph{empty-patch} fields; this is arithmetic over selected
archives, not a rerun of the evaluator. Four of five selected archives agree arithmetically
with a reported point at the stated denominator of 300. \texttt{release-32b-zero} contains 57
resolved and 73 empty-patch tasks against 19.0\% and 24.3\%;
\texttt{release-32b-exp0} contains 59 and 55 against 19.7\% and 18.3\%.

\begin{table}[t]
\caption{Evidence--operation--claim matrix for the released-byte checks. Each row states the
strongest claim carried by that operation, not a study-wide reproduction verdict.}
\label{tab:evidence-operation-claim}
\small
\begin{tabular}{p{0.13\columnwidth}p{0.28\columnwidth}p{0.42\columnwidth}}
\toprule
Study & Operation & Claim ceiling \\
\midrule
SWE-agent & Compare task IDs with a pinned manifest & Task-manifest equality only \\
SWE-Gym & Count selected archives' resolved/empty fields & Selected-subset arithmetic; no evaluator rerun \\
Agentless & Inspect the generation archive & Archive contents only; no score reproduction \\
ACR & Compare run-report counts & Source-unbound candidate arithmetic; no binding or effect \\
\bottomrule
\end{tabular}
\end{table}

The selected SWE-Gym archives also expose two different gaps.

\paragraph{A reported point that no single selected archive contains.}
The 7B middle sequence point requires 27 resolved and 87 empty-patch tasks. The released
archives give 27 with 112, 29 with 87, 23 with 131, and 28 with 92. Each reported component
appears in a different \texttt{top\_k} archive; none contains both, and the selected
candidate archives are exhausted.

\paragraph{Silently dropped run positions.}
Four released archives hold 299 records against a reported denominator of 300, each missing
a different instance: \texttt{sympy\_\_sympy-13773}, \texttt{sympy\_\_sympy-16106},
\texttt{sympy\_\_sympy-14396} and \texttt{sympy\_\_sympy-11897}. Arithmetic under denominator 300 agrees with the displayed proportions,
so the absence is invisible in their rounded presentation. Four different instances across four distinct
public run archives land in one repository, and no terminal reason is recorded for any of them.

\paragraph{Checks with narrower scopes.}
For SWE-agent, the check establishes task-manifest equality with the pinned public revision
and nothing beyond it. For Agentless, we inspect the released generation archive; we do not
claim an evaluator rerun or full figure recomputation for either study.

\paragraph{Three reported repetitions and unbound run-output reports.}
AutoCodeRover~\cite{AutocoderoverAutonomousProgram2024} reports three repetitions of its base
configuration, cites the standard $k=1$ estimator, and annotates two of its three printed Lite
figures accordingly: 19.00\% (57) at @1 and 26.00\% (78) at @3. The third, ACR-sbfl at
22.00\% (66), carries neither annotation, and the source states no run count for it anywhere.
The displayed pair therefore contrasts a three-run statistic with a figure whose aggregation
the source does not state. Its released
run-output reports expose base-run count triples of 54/57/57 under the new evaluation and
48/47/50 under the old. The corresponding ACR-sbfl reports contain 66 and 61. Raw
per-repetition contrasts are therefore \{9,12\} under the new reports and \{11,13,14\} under
the old, spanning 9--14. Applying the cited estimator to each complete triple instead gives
source-unbound full-bundle differences of $66-56=10$ and $61-145/3=38/3$. The source does
not bind a table cell to a public bundle or state the complete selection/count presentation
transform. Neither arithmetic selects an effect or changes the identifiability verdict.

\subsection{Joint identifiability; tier-two compatibility analysis not executed}
\label{sec:results:compatibility}

Each of the 20 comparison decisions nested in the four studies was adjudicated under a
protocol whose rule identity and inputs are digest-recorded. The digests establish which rule
and inputs were used, not that the rule preceded the verdicts: the protocol and first verdicts
entered the repository together. Identification requires that a reproducer can determine the
tuple \emph{(release set, run-output set, evaluator and
configuration, task universe, aggregation rule)} without selecting among candidates by
comparing their outputs to the reported value. The outcome-specific v2 re-adjudication assigns
all 20 decisions \textsc{unidentified}; none is \textsc{identified}, \textsc{open-universe
singleton}, \textsc{no admissible tuple}, or \textsc{compatibility not evaluated}. All five
condition groups across the four studies are consequently \textsc{not identified}. Each
decision lacks at least one enumerable candidate of the type required by its reported outcome.
This is non-exposure in the public record, not proof that the object never existed.

The earlier v1 artifact-global enumeration reported 18 \textsc{unidentified} decisions and
set two aside after constructing an upper bound of 13{,}200 combinations per decision. That
record is preserved as history but superseded: it mixed outcome records with evaluator
procedures, treated generic documents as aggregation rules, and did not represent the
presentation transform. Neither the 18+2 split nor the 13{,}200 upper bound is a current
scientific result.

For CMP-034, the source's three-repetition statement covers the base configuration alone,
and neither it, the cited $k=1$ estimator, nor the printed 57/66/78 counts bind the $k=1$
value to a public run-output bundle. As the arithmetic
in \S\ref{sec:results:recompute} shows, the 9--14 range consists of raw per-repetition
contrasts, while the complete-bundle source-unbound differences are 10 and $38/3$. Any effect
claim remains conditional on the unreported source-to-public-bundle binding and complete
selection/count presentation transform; selecting one would be reverse inference from the
reported value.

For CMP-035, the analysis unit is a unique SWE-bench Lite task. The denominator is the 78
unique tasks with any plausible selected final patch across the three repetitions, and the
numerator is the 51 unique tasks with any semantically correct selected final patch across
those repetitions. The source reports those aggregates,
but the public record contains no task-level human semantic-label vector. Machine evaluator
reports are not candidates for that human outcome, so the outcome-specific run-output
component is unenumerable and the decision is \textsc{unidentified}.

The planned tier-two per-class compatibility analysis was not run. Accordingly, the v2
identifiability result is not accompanied by transfer-attemptability or per-class compatibility
claims.

\subsection{Post-hoc sensitivity of the identifiability result}
\label{sec:results:sensitivity}

The v2 sensitivity analysis was written after the all-negative result was known. R0 is the
recorded criterion; R1 drops candidate-universe closure; R2 additionally drops aggregation;
R3 additionally drops the evaluator; and R4 additionally drops the release set. R5 is a
separate availability-only comparator that restores all five components and drops uniqueness,
not a continuation of R4. No case consults a reported value or supplies missing evidence.
R0 preserves the adjudicated state, and Table~\ref{tab:sensitivity} gives the counts at every
rung. A singleton created only by removing a coordinate in
R2--R4 remains compatibility-not-evaluated unless it descends from an already compatible
singleton; projected compatibility is not invented.

\begin{table}[t]
\caption{Post-hoc v2 sensitivity over the same 20 decisions. NU means that multiple
enumerated candidates remain; it is not an admissibility verdict.}
\label{tab:sensitivity}
\begin{tabular}{lrrr}
\toprule
Case & Unidentified & NU & Identified/available \\
\midrule
R0 & 20 & 0 & 0 \\
R1 & 20 & 0 & 0 \\
R2 & 20 & 0 & 0 \\
R3 & 19 & 1 & 0 \\
R4 & 17 & 3 & 0 \\
R5 & 20 & 0 & 0 \\
\bottomrule
\end{tabular}
\end{table}

No rung yields an already evaluated, supported identification; R3/R4 NU counts are raw
candidate multiplicity, not projected admissibility verdicts. The result is therefore not
created by closure alone,
although the absence of an actual-deposit known-positive control prevents an empirical
sensitivity estimate for the rule itself.

\subsection{Corpus-scale availability}
\label{sec:results:corpus}

This subsection characterises the cohort the four audited studies sit in. It is not a
prevalence estimate: the roster of 47 is a bounded, protocol-defined discovery cohort frozen
before outcome coding, not an enumeration of the eligible literature, and every count below
names its own denominator. Twenty-five of the 47 rostered papers expose an artifact locator;
the remaining 22 are dispositioned in \S\ref{sec:threats}.

Judgement is by reading, not by pattern: for each of those 25 a blind evidence dump of the
deposit's file listing and prose windows was formed with every automated verdict withheld,
and the questions below were answered from it. A pattern screen was run alongside for
comparison and disagreed with the reading on 8 of the 25; \S\ref{sec:threats} reports what
it got wrong and why the reading, not the screen, carries what follows.

The cohort's earlier binding grades are withdrawn. They were formed from prose statements in
the deposits -- a README pairing an expected output file with a reported value, a document
saying its table is the figure in the paper -- and a prose statement is neither necessary nor
sufficient for the property this paper is about. A manifest, an immutable commit and a run
index can bind a released artifact to a reported number with no such sentence, and a sentence
naming a release still leaves the number unbound when the run-output set, evaluator
configuration, task universe or aggregation rule remains ambiguous.

In their place we report what this evidence can carry. The same five components were
adjudicated for each deposit by a second model family, reading the frozen protocol and a blind
evidence dump with the earlier verdicts withheld, and recording for each component whether the
deposit exposes exactly one candidate, several, or none. This is an exposure screen, not an
identifiability verdict: the evidence is listings and prose rather than file contents, so it
cannot establish the uniqueness or the candidate-universe closure the protocol requires, and a
component exposed once here is not thereby identified.

No deposit among the 20 exposes a single candidate for all five components. Sixteen expose
none at all for at least one of them. The component that is unexposed most often is the same
one that is unreached in every deep-audit decision: 15 of the 20 deposits expose no candidate
release or commit for the runs behind their reported figures, and only three expose exactly
one, while task universes and evaluators are exposed several times over in 15 and 13 deposits
respectively. The pattern in the four studies audited at
decision level is not peculiar to them.

Cross-adjudication also corrected this paper. Two studies' dispositions had been recorded
against each other's titles; the second adjudicator disagreed on exactly those two rows and
the evidence settles which way round they go. The aggregate counts are unchanged, because the
error exchanged one \textsc{yes} for one \textsc{no}, but the rows were wrong and are
corrected in the artifact.

The first question is whose artifact the deposit is, and Table~\ref{tab:corpus} gives the
resulting disposition. Five of the 25 locators are not the
paper's own: three are third-party agents or models the paper evaluated against or built on,
one is a GitHub account path naming no repository, and one is a set of applications the paper
scanned. A locator is not a release, and counting these five as artifacts would credit their
studies with evidence belonging to someone else.

\begin{table}[t]
\caption{Disposition of the 25 rostered papers that expose an artifact locator, of which 20
expose one of their own. Counts are the disposition of this cohort, not a rate in the
literature. \textsc{not found} is not absence: the judgement reads the deposit's listing and
prose. \textsc{present elsewhere} is one study whose traces, predictions and harness logs are
published as linked downloads rather than inside the repository.}
\label{tab:corpus}
\begin{tabular}{lr}
\toprule
Disposition among the 25 locator-exposing papers & Papers \\
\midrule
Locator is not the paper's own artifact & 5 \\
Own artifact, run outputs \textsc{present} & 12 \\
Own artifact, run outputs \textsc{present elsewhere} & 1 \\
Own artifact, run outputs \textsc{not found} & 4 \\
Own artifact, \textsc{undetermined} (opaque archive) & 2 \\
Own artifact, deposit \textsc{unreachable} & 1 \\
\bottomrule
\end{tabular}
\end{table}

Among the 20 papers in this cohort that expose an artifact of their own, 13 publish run
outputs. One further deposit is unreachable, two hold only archives whose names decide
nothing, and four publish none. These are counts within the frozen roster and among papers
with completed disposition coding; they are not estimates of how often coding-agent studies
publish run outputs.

Of the 22 rostered papers exposing no locator at all, 10 are determinate -- the full text was
read and names no artifact host -- and 12 are undetermined, 9 because no open full text was
located and 3 because the located text was not retrievable. A paywall is not counted as an
absence.
